% Section 1. Introduction. Carries all contributions (C1-C4) and the not-claimed list.
% Evidence map: paper/review/OUTLINE.md section 2, "1 Introduction".
\section{Introduction}
\label{sec:intro}

A diffusion language model (DLM) generates text by repeatedly predicting many positions in parallel and committing a subset of them, rather than emitting one token at a time from left to right. \citet{arXiv:2107.03006} introduced diffusion on categorical state spaces; later work specialised it to the masked, or absorbing-state, case~\cite{arXiv:2310.16834,arXiv:2406.07524,arXiv:2406.04329,arXiv:2406.03736}, developed a parallel line that diffuses continuous embeddings~\cite{arXiv:2205.14217}, and trained such models at the scale of autoregressive language models~\cite{arXiv:2502.09992}. The same machinery now underlies unified models that understand and generate more than one modality with a single network~\cite{arXiv:2503.20853,arXiv:2505.15809,arXiv:2408.12528}. The literature is growing quickly: of the 540 records with an arXiv identifier in the curated list that seeds our corpus, 243 first appeared in 2026 and only 19 in 2022 or earlier.

This article reviews that literature systematically. From 79,963 identified records, 615 studies were included; a stratified random sample of them was extracted in full and appraised for quality. Beyond describing methods, the review organises the field around five design choices that every generator makes, asks how well the published evidence supports the field's main claims, and records how the headline results of a random sample of papers were decoded. The rest of this section explains why such a review is needed (Section~\ref{sec:intro-why}), positions it against earlier surveys (Section~\ref{sec:intro-related}), and states the research questions, the contributions and the organisation of the article.

\subsection{Decoding choices and the comparability of reported results}
\label{sec:intro-why}


Comparing results across this literature is hard for a reason particular to how these models are used: a DLM has no canonical way to sample from it. The number a paper reports depends on how many positions are committed per step, which positions, and how the committed values are chosen (the \emph{value rule}): drawn at temperature one, sharpened, truncated, guided, or selected from several candidates. Papers differ in all of these and often do not say which they used. Recent evaluation studies have begun to document the consequence. One reports that the ranking of parallel-decoding methods depends on the prompt template~\cite{arXiv:2606.29228}; another examines how efficiency is evaluated across papers~\cite{arXiv:2510.18480}; a technical note discusses benchmark choice and the relation between generative perplexity and entropy~\cite{arXiv:2604.02718}. % src: research/review/survey-landscape.md sections 1c and 2 (E1: abstracts read, not full texts); research/ledgers/papers.jsonl P0079, P0078, P0023, P0113

This review asks how well the published evidence supports the field's claims, and it organises the design choices of these models so that the question can be asked precisely. Search and selection follow PRISMA~2020~\cite{page2021prisma}, and each body of evidence is judged rather than only described.

\subsection{Prior surveys of diffusion language and unified multimodal models}
\label{sec:intro-related}

Two surveys are closest to ours in scope. The survey of \citet{arXiv:2508.10875} (version~3, over 220 references) covers continuous, discrete and hybrid models, training, inference, multimodal models and applications. Its performance study compares DLMs with autoregressive models on eight benchmarks using numbers that are, in its words, primarily taken from the original publications. % src: research/review/survey-landscape.md section 1a (row 2508.10875 v3: refs 228; Sec VI Performance Study)
\citet{arXiv:2506.13759} (version~5, 175 references) cover the discrete case with catalogue tables. % src: research/review/survey-landscape.md section 1a (row 2506.13759 v5: refs 175)
Other surveys address unified multimodal models~\cite{arXiv:2505.02567,arXiv:2409.14993,arXiv:2605.25343}, parallel text generation~\cite{arXiv:2508.08712}, inference efficiency~\cite{arXiv:2607.12829}, edge deployment~\cite{arXiv:2609.04778}, alternatives to next-token prediction~\cite{arXiv:2509.24435} and earlier diffusion-for-text work~\cite{arXiv:2303.06574}.

In a keyword scan of the full texts of sixteen related surveys and tutorials, none states a search protocol with database strings, screening counts and a flow diagram. % src: research/review/survey-landscape.md claims block and section 1a. Reconciled 2026-10-21: the claim block's "14" predates two rows added to table 1a (arXiv:2512.05092, a tutorial, protocol "not applicable"; arXiv:2608.20743, protocol "none found"), so all 16 rows lack a stated search protocol.
One presents itself as a survey conducted through a systematic review and gives venues, years and inclusion criteria, but no search strings or counts~\cite{arXiv:2509.24435}; another gives inclusion criteria for one subsection~\cite{arXiv:2508.08712}. % src: survey-landscape.md section 1a (rows 2509.24435, 2508.08712)
None of them audits what the surveyed papers report: seeds, intervals, matched compute, or the sampling rule behind the headline number. % src: survey-landscape.md section 1a, columns "Audits reporting" (none found in any Tier-1 row)
Neither close survey places models in a design table whose columns are design choices. \citet{arXiv:2506.13759} give no comparison table, and the 36-model table of \citet{arXiv:2508.10875} has mostly metadata columns, of which only the diffusion type is a design choice. % src: research/review/dossier/competitor-anatomy.md section 3, verdict ii
The survey of \citet{arXiv:2508.10875} also states the independence of parallel predictions as a qualitative problem, without a decomposition. % src: competitor-anatomy.md section 3, verdicts iii and iv
We do not claim to be the first survey of this area; the contribution is a different kind of review: a documented search and selection, a design space tested on named models, a census of how results are decoded and reported, and an error budget used to classify methods.

\subsection{Research questions and evidence base}
\label{sec:intro-rq}

The protocol, which follows the guidelines of Kitchenham and Charters for systematic reviews in computing~\cite{kitchenham2007guidelines}, fixes six research questions, each answered from evidence in a dedicated subsection; RQ2 to RQ5 end with a verdict (well supported, weakly supported, mixed or untested). % src: research/review/protocol.md sections 1 and 7
\begin{description}
  \item[RQ1] Which formulations of diffusion language models and unified multimodal diffusion models exist, and what separates them (state space, forward process, objective, parameterisation, sampler)?
  \item[RQ2] How do these models compare with autoregressive models, and under which protocols were the comparisons made?
  \item[RQ3] What do training strategies contribute (pretraining recipe, instruction or reasoning fine-tuning, reinforcement-learning post-training, distillation)?
  \item[RQ4] How does inference trade quality against compute (denoising steps, unmasking order, remasking, caching, parallel decoding)?
  \item[RQ5] In unified multimodal models, what evidence supports, limits or contradicts understanding--generation synergy?
  \item[RQ6] How rigorous is the evidence base (reporting, reproducibility, statistics, agreement between a paper's claims and its own tables)?
\end{description}
% src for the six questions: research/review/protocol.md section 1 (wording shortened only in RQ5 by dropping "evidence used")

Not every question rests on every included study. RQ1 is answered from the screening fields of all 615 included studies, from the 47 models of the design-space table (Section~\ref{sec:design}) and from a stratified sample of the included studies extracted in full. RQ2 to RQ6 are answered from that sample (59 included studies, with 20 seed-corpus studies reported apart), with shares weighted to the size of each stratum and given with intervals (Section~\ref{sec:rq}). The census of Section~\ref{sec:census} reads a random sample of the seed corpus to code the value rule behind each headline number, and the seed-corpus stratum of the extraction sample is drawn from it. % src: sections/05-rq-evidence.tex; analysis/extraction/d0052/ (D0051, D0052, D0054, D0060)

Too few comparisons in the literature share a task, a metric and a compute budget for a meta-analysis, so the review reports counts and weighted shares with intervals rather than a pooled estimate.

\subsection{Contributions of this review}
\label{sec:intro-contrib}

The review makes four contributions and credits the work each builds on. Section~\ref{sec:practice} condenses them for practitioners as canonical models mapped onto the five choices, decoding defaults and a reporting checklist.

\paragraph{A design space in one notation, tested on 47 models.}
Every generator in scope is written as a choice of channel, reveal law, predictor, commitment policy and revision rule, with fourteen levers that change quality or cost (Sections~\ref{sec:notation} and~\ref{sec:design}). The notation builds on published unification results for the masked case~\cite{arXiv:2406.07524,arXiv:2406.04329,arXiv:2406.03736}; we do not claim them. Of the 47 models, 36 fit the notation fully and 11 fit it partly, and every model could be placed; the 11 partial fits mark the stated limits of the notation. % src: research/review/coverage/design-space-v1.csv (column "fits": yes 36, partly 11; counted by script 2026-10-15); design-space-v1.md claims block
The test is less demanding than the count suggests, because the vocabulary was written to be inclusive and 29 of the 47 rows were taken from earlier readings of primary sources rather than re-read for this table. % src: design-space-v1.md "competing_explanations" and "How to read the table" (M01 to M29 copied from the primers)

\paragraph{A census of value rules and reporting.}
In a random sample of 176 records drawn from the 1,106 that both screeners passed to full text and that have an arXiv identifier, 167 papers were eligible. For 153 of them the value rule behind the headline result can be coded (152 report a generated-sample metric and one shows samples only). Of those 153, 147 (96\%, 95\% interval 92 to 98) either decode the headline with a value rule other than plain sampling at temperature one, or do not state the rule. Of the 152 with a metric, 28 pair the headline with a diversity or entropy measure, and of the 55 that report the headline at two or more step settings, 20 state that the value rule was held fixed (Section~\ref{sec:census}). % src: research/review/coverage/census-result-v2.md ("Result", commit 8dd2af3); analysis/census/summary-random.log (final block)
That reported quality and fairness depend on the sampling or evaluation setting has been argued for other metrics and model families~\cite{arXiv:2606.29228,arXiv:2606.20536,arXiv:2607.24763}; the census quantifies the reporting gap for diffusion language models. % src: census-result-v2.md ("What can be said", credit sentence); decisions D0019, D0021

\paragraph{The error budget as a taxonomy.}
For any sampler, described by its commit kernel (the joint law of the positions it commits and the values it gives them; Section~\ref{sec:notation-process}), the divergence between data and sampler along a data path with the same position law splits into a factorisation term, a learning term and one signed value term $V$ that collects every departure of the committed values from the product of the predicted marginals (Section~\ref{sec:budget}). The two-term case $V=0$ is published~\cite{arXiv:2505.24857}; the three-term form is a corollary of the chain rule that we write out. Other pieces are published and we credit them: the majorisation argument for a one-token commit and the sample--select--resample correction~\cite{arXiv:2605.12522}; the distinction between sampling then choosing and choosing then sampling, and the closeness of noised MaskGIT~\cite{arXiv:2202.04200} to a tilted sampler~\cite{arXiv:2510.04525}; the split between confidence-based remasking and top-probability selection~\cite{arXiv:2609.34509}; and the best-of-$N$ KL and win-rate bounds~\cite{arXiv:2401.01879}. % src: research/review/framework-v2.md sections A and B; dossier/prior-art-selection.md
We treat $V$ as a bookkeeping device that locates where a sampler's value rule acts. It is not a quantity with a canonical size: the split depends on the data-path coupling and, for any ``selection'' share, on a reference policy. % src: framework-v2.md section B (paragraph beginning "The decomposition is not unique")

\paragraph{Open problems with the evidence each needs.}
Each open problem is stated with what is known now, the measurement that would settle it, and what would refute it (Section~\ref{sec:open}). They concern what a step sweep identifies and what it does not; whether a revision sampler's gain survives matching the value rule; whether cross-modal dependence in a unified model is concentrated in the early steps, and whether schedules that avoid committing dependent text and image positions together help; and a four-arm measurement on one frozen checkpoint that separates sample-then-filter selection from its resampled and tilted variants. % src: research/review/framework-v2.md sections C (P1, P3, P5) and D item 4

\subsection{What the review does not claim}
\label{sec:intro-notclaimed}

The notation, the census and the error budget serve the survey: they organise and evaluate the published work rather than propose new methods. We do not claim a selection error or a tilt cost as new error terms; we examined both and set them aside. % src: research/ledgers/decisions.jsonl D0011; research/review/framework-v2.md section A
We make no claim of a per-step impossibility result, of the size of the value term, the coupling term or the cross-modal dependence on a real model, or of the U-shaped dependence on the number of steps that appears in a toy model. % src: framework-v2.md section D "Not claimed"
We do not claim that the census sample is a sample of the whole field: it is a random sample of records that passed title-and-abstract screening and have an arXiv identifier. We do not claim that sampler effects on joint text--image consistency are new, or that the value rule carries that consistency; the single preregistered run of Section~\ref{sec:exp003} supports neither claim. % src: research/review/audits/2026-10-15-gate-EXP-003.md section 6 items 4 and 7; decision D0019

\subsection{Organisation of the article}
\label{sec:intro-guide}

\suppFigTree{} in the Supplement places the 47 models of the design space (Section~\ref{sec:design}) in a taxonomy by state space and model family, and gives an entry point to the literature: each leaf cites the model's paper. % src: research/review/coverage/design-space-v1.csv ("family"); scripts/fig_taxonomy.py


Section~\ref{sec:method} gives the review method, Section~\ref{sec:notation} the notation, Section~\ref{sec:design} the design space, Section~\ref{sec:rq} the evidence by research question, Section~\ref{sec:census} the census, Section~\ref{sec:budget} the error budget and Section~\ref{sec:practice} what the evidence means for practice. Section~\ref{sec:open} sets out the open problems, with a preregistered check on a frozen unified model (Section~\ref{sec:exp003}), and Section~\ref{sec:limits} the threats to validity. The full search strings, the PRISMA checklist, the codebooks and the row-level tables are in the Supplement.
